\documentclass{beamer}

\usepackage{../packages/packages_mathalea}
\input{../competences/competences_latex.tex}
\usepackage{fontawesome}
\begin{document}

%%%%%%%%%%%%%%%%%%%%%%%%%%%%%%%%%%%%%%%%%%%%%%%%%%%%%%%%%%%%%%%%%%%%%%%%%%%%%%%
%%%%%%%%%%%%%%%%%%%%%%%%%%%%%%%%%%%%%%%%%%%%%%%%%%%%%%%%%%%%%%%%%%%%%%%%%%%%%%%
%%%%%%%%%%%%%%%%%%%%%%%%%%%%%%%%%%%%%%%%%%%%%%%%%%%%%%%%%%%%%%%%%%%%%%%%%%%%%%%

\begin{frame}{\GetCompetence{Fonc4}}
	% increment the counter
	\onslide<1->{ Déterminer les zéros de la fonction $f(x)=x^2-5x+6$.}

	\onslide<2->{{\bfseries Solution~:} Déterminer les zéros de $f$ revient à déterminer la préimage de 0 par $f$, c'est-à-dire à résoudre l'équation \[x^2-5x+6=0.\] 

	Les solutions sont $S=\left\{2;3\right\}$, donc $Z_f=\left\{2;3\right\}$.}
\end{frame}

\end{document}
